\documentclass[10pt,a4paper]{amsart}
\usepackage{amsmath,amssymb,amsthm,graphicx}
% Retained minimal preamble for the Stacks Project chapter varieties.tex.
% Verbatim excerpt of the chapter preamble: Xy-pic setup (native preamble lines 14-18), the native theorem
% declarations with the mathematical macros (native preamble lines 154-217), and the project own no-op
% definitions of the reference/slogan/history display environments (native wrapper.tex lines 1-3) that the
% counted statement uses for its slogan. Copied mechanically; nothing added, removed or redefined.
% For commutative diagrams we use Xy-pic
\usepackage[all]{xy}

% We use 2cell for 2-commutative diagrams.
\xyoption{2cell}
% Theorem environments.
%
\theoremstyle{plain}
\newtheorem{theorem}[subsection]{Theorem}
\newtheorem{proposition}[subsection]{Proposition}
\newtheorem{lemma}[subsection]{Lemma}

\theoremstyle{definition}
\newtheorem{definition}[subsection]{Definition}
\newtheorem{example}[subsection]{Example}
\newtheorem{exercise}[subsection]{Exercise}
\newtheorem{situation}[subsection]{Situation}

\theoremstyle{remark}
\newtheorem{remark}[subsection]{Remark}
\newtheorem{remarks}[subsection]{Remarks}

\numberwithin{equation}{subsection}

% Macros
%
\def\lim{\mathop{\mathrm{lim}}\nolimits}
\def\colim{\mathop{\mathrm{colim}}\nolimits}
\def\Spec{\mathop{\mathrm{Spec}}}
\def\Hom{\mathop{\mathrm{Hom}}\nolimits}
\def\Ext{\mathop{\mathrm{Ext}}\nolimits}
\def\SheafHom{\mathop{\mathcal{H}\!\mathit{om}}\nolimits}
\def\SheafExt{\mathop{\mathcal{E}\!\mathit{xt}}\nolimits}
\def\Sch{\mathit{Sch}}
\def\Mor{\mathop{\mathrm{Mor}}\nolimits}
\def\Ob{\mathop{\mathrm{Ob}}\nolimits}
\def\Sh{\mathop{\mathit{Sh}}\nolimits}
\def\NL{\mathop{N\!L}\nolimits}
\def\CH{\mathop{\mathrm{CH}}\nolimits}
\def\proetale{{pro\text{-}\acute{e}tale}}
\def\etale{{\acute{e}tale}}
\def\QCoh{\mathit{QCoh}}
\def\Ker{\mathop{\mathrm{Ker}}}
\def\Im{\mathop{\mathrm{Im}}}
\def\Coker{\mathop{\mathrm{Coker}}}
\def\Coim{\mathop{\mathrm{Coim}}}

% Boxtimes
%
\DeclareMathSymbol{\boxtimes}{\mathbin}{AMSa}{"02}

%
% Macros for moduli stacks/spaces
%
\def\QCohstack{\mathcal{QC}\!\mathit{oh}}
\def\Cohstack{\mathcal{C}\!\mathit{oh}}
\def\Spacesstack{\mathcal{S}\!\mathit{paces}}
\def\Quotfunctor{\mathrm{Quot}}
\def\Hilbfunctor{\mathrm{Hilb}}
\def\Curvesstack{\mathcal{C}\!\mathit{urves}}
\def\Polarizedstack{\mathcal{P}\!\mathit{olarized}}
\def\Complexesstack{\mathcal{C}\!\mathit{omplexes}}
% \Pic is the operator that assigns to X its picard group, usage \Pic(X)
% \Picardstack_{X/B} denotes the Picard stack of X over B
% \Picardfunctor_{X/B} denotes the Picard functor of X over B
\def\Pic{\mathop{\mathrm{Pic}}\nolimits}
\def\Picardstack{\mathcal{P}\!\mathit{ic}}
\def\Picardfunctor{\mathrm{Pic}}
\def\Deformationcategory{\mathcal{D}\!\mathit{ef}}
\newenvironment{reference}{}{}
\newenvironment{slogan}{}{}
\newenvironment{history}{}{}
\usepackage{cleveref}
\crefname{originaldefinition}{original definition}{original definitions}
\Crefname{originaldefinition}{Original definition}{Original definitions}
\makeatletter
\newlabel{algebra-definition-geometrically-normal}{{0380}{source}{}{sourcepointer-1}{}}
\newlabel{algebra-definition-geometrically-normal@cref}{{[originaldefinition][0][]0380}{[1][0][]source}{}{}{}}
\makeatother
\crefname{originalsection}{original section}{original sections}
\Crefname{originalsection}{Original section}{Original sections}
\makeatletter
\newlabel{schemes-section-fibre-products}{{01JO}{source}{}{sourcepointer-2}{}}
\newlabel{schemes-section-fibre-products@cref}{{[originalsection][0][]01JO}{[1][0][]source}{}{}{}}
\makeatother
\title{Six equivalent characterizations of geometric normality for a scheme over a field}
\author{The Stacks Project Authors}
\date{stacks-project:varieties.tex:local-snapshot-20261011; GNU FDL 1.2 or later}
\begin{document}
\maketitle
\noindent\textit{Editorial note.} All delivered mathematical text is the author's own, copied from the pinned Stacks Project chapter file varieties.tex; no statement or proof step was rewritten, reordered or removed.\par\medskip
\noindent\textit{Editorial note.} The proof is complete as written and closes the cycle of implications; no step is referred to the reader.\par\medskip
\noindent\textit{Editorial note.} The counted spans carry 2 navigation stubs and no bibliography citation.\par\medskip
\noindent\textit{Editorial note.} The preamble is the chapter's retained minimal preamble (sources/stacks-varieties/retained-preamble.tex lines 1-77), a verbatim excerpt of the native preamble.\par\medskip
\noindent\textit{Editorial note.} The label lemma-geometrically-normal is defined only in varieties.tex, so no cross-chapter restatement had to be weighed for this card.\par\medskip
\noindent\textit{Allowed standard prerequisites (editorial).}\begin{itemize}
\item Geometric normality of a k-algebra. Algebra, Definition algebra-definition-geometrically-normal, invoked at proof lines 2192-2193 and named in the statement. Supplies the algebraic property to which conditions (4) and (5) are translated.
\item Base change of schemes as gluing of spectra of tensor products of coordinate rings. Schemes, Section schemes-section-fibre-products, invoked at proof lines 2197-2200. Supplies that normality of X\_\{k'\} follows from normality of the tensor products of the coordinate rings of affine opens.
\end{itemize}\medskip
\ifdefined\hypertarget\hypertarget{sourcepointer-1}{}\fi
\noindent\textit{Source reference.} Original definition 0380: sources/stacks-chow/tags.txt, source lines 4219--4219. Pointer outside the counted span, invoked once for context; the load-bearing argument of the counted statement is the author's proof reproduced in full in the proof span.\par\medskip
\ifdefined\hypertarget\hypertarget{sourcepointer-2}{}\fi
\noindent\textit{Source reference.} Original section 01JO: sources/stacks-chow/tags.txt, source lines 2021--2021. Pointer outside the counted span, invoked once for context; the load-bearing argument of the counted statement is the author's proof reproduced in full in the proof span.\par\medskip
\begin{lemma}
\label{lemma-geometrically-normal}
Let $k$ be a field.
Let $X$ be a scheme over $k$.
The following are equivalent
\begin{enumerate}
\item $X$ is geometrically normal,
\item $X_{k'}$ is a normal scheme for every field extension $k'/k$,
\item $X_{k'}$ is a normal scheme for every finitely generated field
extension $k'/k$,
\item $X_{k'}$ is a normal scheme for every finite purely inseparable
field extension $k'/k$,
\item for every affine open $U \subset X$ the ring $\mathcal{O}_X(U)$
is geometrically normal (see
Algebra, Definition \ref{algebra-definition-geometrically-normal}), and
\item $X_{k^{perf}}$ is a normal scheme.
\end{enumerate}
\end{lemma}

\begin{proof}
Assume (1). Then for every field extension $k'/k$ and
every point $x' \in X_{k'}$ the local ring of $X_{k'}$ at $x'$
is normal. By definition this means that $X_{k'}$ is normal.
Hence (2).

\medskip\noindent
It is clear that (2) implies (3) implies (4).

\medskip\noindent
Assume (4) and let $U \subset X$ be an affine open subscheme.
Then $U_{k'}$ is a normal scheme for any finite purely inseparable
extension $k'/k$ (including $k = k'$). This means that
$k' \otimes_k \mathcal{O}(U)$ is a normal ring for all
finite purely inseparable extensions $k'/k$. Hence
$\mathcal{O}(U)$ is a geometrically normal $k$-algebra by definition.
Hence (4) implies (5).

\medskip\noindent
Assume (5). For any field extension $k'/k$ the base
change $X_{k'}$ is gotten by gluing the spectra of the
rings $\mathcal{O}_X(U) \otimes_k k'$ where $U$ is affine open
in $X$ (see Schemes, Section \ref{schemes-section-fibre-products}).
Hence $X_{k'}$ is normal. So (1) holds.

\medskip\noindent
The equivalence of (5) and (6) follows from the definition
of geometrically normal algebras and the equivalence (just proved)
of (3) and (4).
\end{proof}

\end{document}
